\documentclass[10pt,a4paper]{amsart}
\usepackage[margin=19mm]{geometry}
\usepackage{amsmath,amssymb,mathtools}
\usepackage[scr=rsfs,cal=euler]{mathalfa}
\usepackage{xfrac}
\usepackage[unicode,hidelinks,bookmarks=false]{hyperref}
\newcommand{\ie}{i.e.\ }
\newcommand{\cf}{cf.\ }
\newcommand{\resp}{respectively\ }
\newcommand{\Subsection}{\subsection}
\providecommand{\nobreakdash}{\nobreak\hbox{-}}
\setlength{\emergencystretch}{2em}
\multlinegap0pt
\usepackage{bbm}
\usepackage{graphicx}
\numberwithin{equation}{section}
\newcommand{\noopsort}[1]{}
\newcommand{\id}{id}
\newcommand{\N}{\mathbbm{N}}
\newcommand{\Z}{\mathbbm{Z}}
\newcommand{\Q}{\mathbbm{Q}}
\newcommand{\R}{\mathbbm{R}}
\newcommand{\D}{\mathbb{D}}
\newcommand{\Hy}{\mathbbm{H}}



\newcommand{\Ham}{\mathrm{Ham}}
\newcommand{\topo}{\mathrm{top}}
\newcommand{\hofer}{\mathrm{Hofer}}
\newcommand{\Spec}{\mathrm{Spec}}
\newcommand{\Braid}{\mathrm{Braid}}
\newcommand{\rid}{\mathrm{id}}
\newcommand{\Cyl}{\mathrm{Cyl}}
\newcommand{\Cross}{\mathrm{Cross}}
\newcommand{\rmD}{\mathrm{D}}
\newcommand{\rCZ}{\mathrm{CZ}}
\newcommand{\rtriv}{\mathrm{triv}}
\newcommand{\pr}{\mathrm{pr}}
\newcommand{\intt}{\mathrm{int}}
\newcommand{\diam}{\mathrm{diam}}

\newcommand{\clP}{\mathcal{P}}
\newcommand{\clY}{\mathcal{Y}}
\newcommand{\clA}{\mathcal{A}}
\newcommand{\clI}{\mathcal{I}}
\newcommand{\clF}{\mathcal{F}}
\newcommand{\clU}{\mathcal{U}}
\newcommand{\clB}{\mathcal{B}}
\newcommand{\clQ}{\mathcal{Q}}
\newcommand{\clM}{\mathcal{M}}
\newcommand{\clH}{\mathcal{H}}
\newcommand{\clN}{\mathcal{N}}
\newcommand{\clE}{\mathcal{E}}
\newcommand{\clD}{\mathcal{D}}
\newcommand{\clL}{\mathcal{L}}
\newcommand{\clT}{\mathcal{T}}

\newcommand{\fkb}{\mathfrak{b}}
\newcommand{\fkg}{\mathfrak{g}}
\newcommand{\fkF}{\mathfrak{F}}
\newcommand{\fkG}{\mathfrak{G}}
\newcommand{\fkf}{\mathfrak{f}}
\newcommand{\fki}{\mathfrak{i}}
\newcommand{\fkL}{\mathfrak{L}}
\newcommand{\fkN}{\mathfrak{N}}

\newcommand{\bfu}{\mathbf{u}}
\newcounter{stepcount}
\newenvironment{step}{\refstepcounter{stepcount}\begin{proof}[Step~\thestepcount]}{\let\qed\relax\end{proof}}
\newcommand*{\mk}{\mkern -1mu}
\newcommand*{\Mk}{\mkern -2mu}
\newcommand*{\mK}{\mkern 1mu}
\newcommand*{\MK}{\mkern 2mu}
\newcommand*{\romanenumi}{\renewcommand*{\theenumi}{\roman{enumi}}}
\newcommand*{\Romanenumi}{\renewcommand*{\theenumi}{\Roman{enumi}}}
\newcommand*{\alphenumi}{\renewcommand*{\theenumi}{\alph{enumi}}}
\newcommand*{\Alphenumi}{\renewcommand*{\theenumi}{\Alph{enumi}}}
\let\oldtilde\tilde
\renewcommand*{\tilde}[1]{\mathchoice{\widetilde{#1}}{\widetilde{#1}}{\oldtilde{#1}}{\oldtilde{#1}}}
\let\oldhat\hat
\renewcommand*{\hat}[1]{\mathchoice{\widehat{#1}}{\widehat{#1}}{\oldhat{#1}}{\oldhat{#1}}}
\newtheorem{theo}{Theorem}[section]
\newtheorem{prop}[theo]{Proposition}
\newtheorem{lemm}[theo]{Lemma}
\newtheorem{coro}[theo]{Corollary}
\theoremstyle{definition}
\newtheorem{defi}[theo]{Definition}
\theoremstyle{remark}
\newtheorem{rema}[theo]{Remark}
\title[Entropy approximated by punctured fundamental-group growth]{Positive entropy of a compact-surface diffeomorphism is approximated by fundamental-group growth relative to a single hyperbolic periodic orbit}
\author{Marcelo R. R. Alves and Matthias Meiwes}
\date{Source: Annales Henri Lebesgue \textbf{7} (2024), 521--579; DOI 10.5802/ahl.205; official published version}
\begin{document}
\maketitle
\noindent\textit{Editorial scope.} The counted unit is original Theorem B.1 with its complete authored proof, which includes the deferred complete proof of Proposition B.5 where the punctured fundamental-group growth is actually produced. Retained uncounted core: the publisher-native definition of the growth rate $\Gamma_{\pi_1}([\varphi,\mathcal{P}])$, the Pesin-theoretic Lyapunov-chart and admissible-rectangle paragraph, and Propositions B.2 and B.3 of Katok--Mendoza with the combinatorial Lemma B.4, all copied exactly from the source. The published appendix numbering is reproduced.\par
\noindent\textit{Reference handling.} The definition of $\Gamma_{\pi_1}$ lives in Section 1.4 of the published article, which is not part of this unit; the single cross-reference to it is delivered as a hyperlink to the published article instead of a local label, so the required context is stated in place rather than imported.
\appendix
\setcounter{section}{1}
\section{Approximation of entropy by the entropy of braid types} \label{appendix:approx}

Let $\varphi: \Sigma \to \Sigma$ be a diffeomorphism on a compact surface $\Sigma$, $\clP$ a periodic orbit and $\overline{P}\subset \Sigma$ the associated set of periodic points. Then $\Gamma_{\pi_1}([\varphi, \clP])$ denotes the growth rate of the induced action of $\varphi$ on the fundamental group of $\Sigma \setminus \overline{\clP}$, see Section~\href{https://ahl.centre-mersenne.org/item/AHL_2024__7__521_0/}{Section~1.4}. Also recall that $\Gamma_{\pi_1}([\varphi, \clP]) \leq h_{\topo}(\varphi)$.

The aim of this section is to give a proof of the following result.

\begin{theo}\label{thm:approximation}
Let $\Sigma$ be a compact surface, and let $\varphi: \Sigma\to \Sigma$ be a diffeomorphism such that $h:=h_{\topo}(\varphi)>0$. Then, for any $\epsilon>0$ there is a hyperbolic periodic orbit $\clP$ of $\varphi$ such that $\Gamma_{\pi_1}([\varphi,\clP])> h - \epsilon$.
\end{theo}



This is a variant of a celebrated result of Katok first announced in~\cite{Katok1984}, see also~\cite[Supplement]{Hasselblatt-Katok}, about the approximation of the topological entropy by the entropies on locally maximal hyperbolic invariant sets on which the dynamics is conjugated to a subshift of finite type, and in fact the orbits $\clP$ of Theorem~\ref{thm:approximation} can be found inside those sets. To our knowledge there is no proof of Theorem~\ref{thm:approximation} in the literature, so we include it here. In~\cite{FranksHandel1988} it was proved that there is a collection of orbits for which the exponential growth of $\varphi$ relative to it is positive. We apply the strategy in~\cite{FranksHandel1988} to the hyperbolic horseshoes that arise from Katok and Mendoza's results in~\cite[Supplement]{Hasselblatt-Katok}.

We recall relevant results from~\cite[Suppl.]{Hasselblatt-Katok} which are applications of the theory of non-uniformly hyperbolic dynamics, or Pesin theory, and we refer also to~\cite{BarreiraPesin} and references therein. Consider a $\varphi$-invariant ergodic Borel probability measure $\mu$ on $\Sigma$ and assume it is hyperbolic, i.e. its Lyapunov exponents satisfy $\chi_1< 0< \chi_2$. See~\cite[Suppl., Section~2]{Hasselblatt-Katok} for the definition and existence of Lyapunov exponents of the measure $\mu$, and important properties. It is shown that for almost every $x \in \Sigma$ there are so-called Lyapunov charts $\psi_x: B(0,r(x)) \subset \R^2 \to \Sigma$ for $x$, $r(x)>0$, with respect to which the dynamics has uniformly hyperbolic behaviour. This can be expressed with the notion of admissible stable and unstable manifolds, defined locally in these charts, and which are locally preserved by $\varphi$. In order to do that one restricts $\psi_x$ to $[-\kappa,\kappa]^2$ for sufficiently small\footnote{The number $\kappa>0$ depends on the point $x$, but we will omit this from now on to simplify the notation.} $\kappa>0$ and considers the \emph{regular rectangles} $R(x,\kappa) = \psi_x([-\kappa,\kappa]^2)$ in $\Sigma$. We denote the left boundary of $R$ by $\partial_l R$ given by $\partial_l R = \{ \psi_x(-\kappa, s) \, | \, s \in [-\kappa, \kappa]\}$ and denote the right, bottom and top boundaries of $R= R(x,\kappa)$ by $\partial_r R, \partial_b R$, and $\partial_t R$, respectively, which are defined analogously. Let $\gamma \in (0,1)$. A submanifold $W \subset R(x,\kappa)$ is called an \emph{admissible stable $(\gamma,\kappa)$-manifold near x} if $W= \psi_x(\{(\phi(v),v)\, |\, v\in [-\kappa,\kappa]\})$, where $\phi:[-\kappa,\kappa] \to [-\kappa,\kappa]$ is a $C^1$-map with $\phi(0) \leq \kappa/4$ and $|D\phi| \leq \gamma$. Analogously $W \subset R(x,\kappa)$ is called an \emph{admissible unstable $(\gamma,\kappa)$-manifold near x} if $W= \psi_x(\{(u,\phi(u))\, |\, v\in [-\kappa,\kappa]\}$, where $\phi:[-\kappa,\kappa] \to [-\kappa,\kappa]$ is a $C^1$-map with $\phi(0) \leq \kappa/4$ and $|D\phi| \leq \gamma$. Admissible stable and unstable $(\gamma,\kappa)$-manifolds near $x$ intersect in exactly one point in $R(x,\kappa)$ with angle bounded away from zero and so they endow $R(x,\kappa)$ with a product structure. Moreover one defines admissible unstable (stable) rectangles as sets bounded by two admissible unstable (stable) manifolds. That is, an admissible stable $(\gamma,\kappa)$-rectangle in $R(x,\kappa)$ is a set of the form $V=\psi_x(V')$, $V'= \{(u,v) \in [-\kappa,\kappa]^2 \, | \, u = (1-\tau) \phi_1(v) + \tau \phi_2(v), \, \tau \in [0,1]\}$, where the left boundary $\partial_l V := \psi_x(\{(\phi_1(v),v)\, |\, v\in [-\kappa,\kappa]\}$ and the right boundary $\partial_r V:=$ $\psi_x(\{(\phi_2(v),v)\, |\, v\in [-\kappa,\kappa]\}$ are two admissible unstable $(\gamma,\kappa)$-manifolds near $x$ for which $\phi_1(v) < \phi_2(v)$ for all $v\in [-\kappa, \kappa]$. Analogously, define admissible stable $(\gamma,\kappa)$-rectangles $H$ in $R(x,\kappa)$ with bottom (top) boundaries $\partial_b H$ ($\partial_t H)$. We define $\partial_b V = \partial_b R(x,\kappa) \cap V$, $\partial_t V = \partial_t R(x,\kappa) \cap V$, and similarly $\partial_l H$ and $\partial_r H$. The following statement asserts the existence of rectangular covers and hyperbolic properties of $\phi$ on those rectangles.

\begin{prop}[{\cite[Theorem S.4.16]{Hasselblatt-Katok}}]\label{prop:KatokRectangles}
For every $\delta>0$ and $\rho>0$ there is a compact set $\Lambda_{\delta}$ with $\mu(\Lambda_{\delta})>1-\delta$, and constants $\beta>0$, $\kappa>0$, $\gamma\in (0,1)$ and regular rectangles $R(x_1) = R(x_1,\kappa), R(x_2) = R(x_2,\kappa), \,\ldots,\, R(x_t)= R(x_t,\kappa)$, $t\in \N$, for some $x_1, \,\ldots,\, x_t \in \Lambda_{\delta}$, such that
\begin{enumerate}
\item\label{propB.2.1} $\Lambda_{\delta} \subset \bigcup_{i=1}^t B(x_i, \beta)$ with $B(x_i, \beta) \subset {\intt}R(x_i)$
\item \label{propB.2.2}${\diam}R(x_i) \leq \rho/3$ for $i=1, \,\ldots,\, t$
\item\label{propB.2.3} If $y\in \Lambda_{\delta}, \varphi^n(y) \in \Lambda_{\delta}$ for some $n>0, y \in B(x_i, \beta)$, and $f^n(y) \in B(x_j, \beta)$, then the connected component $V$ of $R(x_i)\cap f^{-n}(R(x_j))$ containing $y$ is an admissible stable $(\gamma, \kappa)$-rectangle near $x_i$ and $H= f^n(V)$ is an admissible unstable $(\gamma, \kappa)$-rectangle near $x_j$.
\item\label{propB.2.4} $\diam f^k(V) < \rho$ for all $0\leq k \leq n$.
\end{enumerate}
\end{prop}

In~\eqref{propB.2.3} above we have that the union $(\partial_l V \cup \partial_r V)$ map to the union of $(\partial_l H \cup \partial_r H)$, etc.

Assume that $h_{\mu}(\varphi)>0$. One can show now the following, see~\cite[Theorem~S.5.9]{Hasselblatt-Katok} and its proof.

\begin{prop}\label{prop:KatokHorseshoe}
Let $\delta>0, \rho>0$ and choose the constants $\beta, \kappa, \gamma$, the set $\Lambda_{\delta}$ and rectangles $R(x_1),\, \ldots,\, R(x_t)$ as in Proposition~\ref{prop:KatokRectangles}. Then for any $\epsilon_1>0$ sufficiently small there is an unbounded set $\clN \subset \N$ such that for all $n\in \clN$ there is $x \in \{x_1, \,\ldots,\, x_t\}$, and $M > e^{n(h_{\mu}(\varphi) - \epsilon_1)}$ disjoint admissible stable $(\gamma,\kappa)$-rectangles $V_1, \,\ldots,\, V_M$ near $x$ that are mapped as in $(3)$ above by $\varphi^n$ to $M$ admissible unstable $(\gamma,\kappa)$-rectangles $H_1, \ldots, H_M$ near $x$.

Moreover,
\[
\Lambda_M =\bigcap_{k\, \in\, \Z} \left(\varphi^{kn}(V_1) \cup \cdots \cup \varphi^{kn}(V_{M})\right)
\]
is a locally maximal hyperbolic invariant set with respect to $\varphi^n$, and $\varphi^n|_{\Lambda_M}$ is topologically conjugate to a full two-sided shift in the symbolic space of $M$ symbols. In particular, $h_{\topo}(\varphi^n|_{\Lambda_M}) > n(h_{\mu}(\varphi) - \epsilon_1)$.
\end{prop}

Note that it follows from the variational principle, the ergodic decomposition theorem and Ruelle's inequality that one can approximate $h_{\topo}(\varphi)$ by $h_{\mu}(\varphi)$ of such measures $\mu$ as considered above, and the conclusions of Proposition~\ref{prop:KatokHorseshoe} imply approximation of the topological entropy $h_{\topo}(\varphi)$ by the topological entropy of hyperbolic horseshoes of $\varphi$.


The proof of Theorem~\ref{thm:approximation} will use the above results and a variant of arguments of Franks and Handel in~\cite{FranksHandel1988}.

\begin{proof}[Proof of Theorem~\ref{thm:approximation}]
Let $\varepsilon>0$ and let $\mu$ be an invariant ergodic hyperbolic measure with $h_{\mu}(\varphi) > h - \epsilon/3$. Choose $\varepsilon_1 = \varepsilon/3$. Let now $\clN\subset \N$ as in Proposition~\ref{prop:KatokHorseshoe}. Choose some $n\in \clN$ with $\frac{ \log(20)}{n} < \varepsilon_1$. Let $x \in \{x_1,\, \ldots,\, x_t\}$ and $M$, $V_1,\, \ldots,\, V_M$, $H_1,\, \ldots,\, H_M \subset R(x)$ as in Proposition~\ref{prop:KatokHorseshoe}. We say that $H_i$ is positively oriented if $\varphi^n(\partial_t V_i) = \partial_t H_i$ and negatively oriented if $\varphi^n(\partial_t V_i) = \partial_b H_i$. We say that $H_i$ lies below $H_j$ in $R(x)$, $i\neq j$, if every path in $R(x)$ from $\partial_b R(x)$ to $\partial_t R(x)$ intersects first $H_i$.


For convenience we restrict to a subset $V'_1,\, \ldots,\, V'_m$ of the admissible stable rectangles $V_1, \,\ldots,\, V_M$ as well as a subset $H'_1, \ldots, H'_m$ of the admissible unstable rectangles \linebreak $H_1,\, \ldots,\, H_M$, both of size $m\geq M/10$ such that
\begin{itemize}
\item $\varphi^n(V'_i) = H'_i$ for $i=1, \,\ldots,\, m$.
\item All rectangles $H'_1,\, \ldots,\, H'_m$ are either all positively oriented or all negatively oriented.
\item Either $H'_1$ is below $H'_2,\, \ldots,\, H'_{m}$ and $H'_m$ is above $H'_1,\, \ldots,\, H'_{m-1}$, or $H'_1$ is above $H'_2, \,\ldots,\, H'_{m}$ and $H'_m$ is below $H'_1, \,\ldots,\, H'_{m-1}$.
\end{itemize}

That we can ensure the last condition, follows from the following observation.

\begin{lemm}
Let $x_1, \,\ldots,\, x_n$ be a sequence of $n$ pairwise different natural numbers. Then it has a subsequence of length $\geq \frac{n}{5}$ such that all elements of the subsequence lie in the closed interval $I$ whose boundary is given by the first and the last element of that subsequence.
\end{lemm}

\begin{proof}
Let $i_0, i_1 \in \{1, \,\ldots,\, n\}$ such that $x_{i_0} = \min \{ x_i \, | \, i = 1,\, \ldots,\, n \}$ and $x_{i_1} = \max \{ x_i \, | \, i = 1,\, \ldots,\, n\}$. We assume that $i_0 < i_1$, the other case is analogous. If $i_1 -i_0 + 1 \geq \frac{n}{5}$, then the subsequence $x_{i_0}, x_{i_0 + 1}, \,\ldots,\, x_{i_1}$ satisfies the properties claimed in the lemma. Otherwise, $i_0> \frac{2n}{5}$ or $n-i_1 > \frac{2n}{5}$. Note that any element of the sequence lies in the interval $[x_1,x_{i_1}]$ or in the interval $[x_{i_0}, x_1]$. So if $i_0>\frac{2n}{5}$, $[x_1,x_{i_1}]$ or $[x_{i_0}, x_1]$ contains $> \frac{n}{5}$ elements from $\{x_1, \,\ldots,\, x_{i_0}\}$. Similarly, if $n-i_1 > \frac{2n}{5}$, the interval $[x_n, x_{i_1}]$ or the interval $[x_{i_0}, x_n]$ contains $> \frac{n}{5}$ elements from $\{x_{i_1},\,\ldots,\, x_{n}\}$. In all situations these elements form a subsequence as claimed.
\end{proof}

We now proceed with the argument for the case that all $H'_i$ are positively oriented, $H'_1$ is below $H'_2, \,\ldots,\, H'_{m}$ and $H'_m$ is above $H'_1, \,\ldots,\, H'_{m-1}$. The other cases are treated analogously. By abuse of notation we drop the symbol $'$, and denote these rectangles by $V_1, \,\ldots,\, V_m$ and $H_1,\, \ldots,\, H_m$. Let $\Lambda_n = \bigcap_{k\in \Z} (\varphi^{kn}(V_1) \cup \cdots \cup \varphi^{kn}(V_{m}))$, and $\Sigma_m$ the set of bi-infinite sequences of $m$ symbols. The map $\theta: \Lambda_n \to \Sigma_m$ defined as $x \mapsto (\theta(x)_k)_{k\, \in\, \Z}$, with $\theta_k(x) = j$ if $\varphi^{kn}(x) \in V_j$, provides a conjugation of $\varphi^n$ with the two-sided shift $\sigma_m$ on $\Sigma_m$.

For any $a_0,\, \ldots,\, a_{k-1} \in \{1, \,\ldots,\, m\}$ we denote by $(a_0\cdots a_{k-1})^{\infty}$ the  $k$-periodic orbit of $\sigma_n$ in $\Sigma_m$ that is given by an infinite repetition of $a_0\cdots a_{k-1}$ in positive and negative direction. If we want to refer to a periodic point of that orbit, we say that $(a_i a_{i+1} \cdots a_{k-1} a_0 \cdots a_{i-1})^{\infty}$ is the point that has $a_i$ at the $0$th position. By abuse of notation we denote the periodic orbits or points of $\varphi^n$ in $\Lambda_n$ that correspond to those in $\Sigma_m$ via $\theta$ also by the symbols $(a_0\cdots a_{k-1})^{\infty}$ etc.

Consider the periodic orbit $\clQ$ of $\varphi^n$ given by
\begin{align*}
\clQ = \left(\prod_{j=2}^{m-1}(mj1j1jmj)\right)^{\infty}.
\end{align*}

We write also $\clQ$ as the orbit
\begin{equation}\label{qqqq}
\left(q_1^2, q_2^2, \ldots, q_8^2, q_1^3, q_2^3 \ldots,\, q_8^3, \,\ldots,\, q_8^{m-2}, q_1^{m-1}, q_2^{m-1},\, \ldots,\, q_8^{m-1}\right),
\end{equation}
where the periodic point $q_1^2$ is given in the symbolic expression by
\[
q_1^2 = \left(\left(\prod_{j=2}^{m-2}(1j1jmjm(j+1))\right) 1(m-1)1(m-1)m(m-1)m2\right)^{\infty},
\]
and the other periodic points $q_l^j$, $j\in \{2,\,\ldots,\, m-1\}$, $l\in \{1,\,\ldots,\,8\}$, accordingly via cyclic permutations of the symbols.
We note the following, which will be relevant below: For any $j \in \{2,\,\ldots,\,m-1\}$
\begin{align}\label{1jm}
\begin{split} 
&q_1^j \in V_1 \cap H_j \cap \varphi^n(H_m), \quad
q_3^j\in V_1 \cap H_j \cap \varphi^n(H_1) \\
& q_5^j\in V_m \cap H_j \cap \varphi^n(H_1), \quad q_7^j\in V_m \cap H_j \cap \varphi^n(H_m),
\end{split}
\end{align}
and
\begin{equation}\label{11mm}
q_2^j,q_4^j, q_6^j, q_8^j \in H_1 \cup H_m.
\end{equation}



We show below that

\begin{prop}\label{prop:entropybound}
$\Gamma_{\pi_1}([\varphi^n, \clQ]) \geq \log(m-2)$.
\end{prop}

From the Proposition we obtain
\begin{align}
\begin{split}
\Gamma_{\pi_1}([\varphi^n, \clQ]) &\geq \log(m-2) \geq \log(M/20) \\ &\geq \log\left(\frac{e^{n(h-2\epsilon_1)}}{20}\right) \geq n(h-2\epsilon_1) - \log(20) \\ &> n(h-\varepsilon).
\end{split}
\end{align}
The hyperbolic periodic orbit $\clQ$ of $\varphi^n$ of period $8(m-2)$ defines a hyperbolic periodic orbit $\clP$ of $\varphi$ of period $8n(m-2)$, with $\overline{\clQ} \subset \overline{\clP}$ for the associated sets of $\clQ$ and~$\clP$. Then $\Gamma_{\pi_1}([\varphi,\overline{\clP}]) = \frac{1}{n}\Gamma_{\pi_1}([\varphi^n,\overline{\clP}]) \geq \frac{1}{n} \Gamma_{\pi_1}([\varphi^n, \overline{\clQ}]) > h-\varepsilon$, which will finish the proof of Theorem~\ref{thm:approximation}.
\end{proof}


\begin{proof}[Proof of Proposition~\ref{prop:entropybound}]
The proof is a variation of an argument from~\cite{FranksHandel1988}.

Identify the universal covering of $\Sigma \setminus \overline{\clQ}$ with the Poincaré disk $\Hy$. This gives us a choice of hyperbolic metric on $\Sigma \setminus \overline{\clQ}$. Every proper arc in $\Sigma \setminus \overline{\clQ}$, i.e. an arc $\alpha$ that lies in $\Sigma\setminus \overline{\clQ}$ up to its endpoints and whose endpoints lie in $\overline{\clQ}$, defines uniquely a geodesic that traces a proper arc homotopic to $\alpha$. Also, any closed curve $\gamma$ in $\Sigma\setminus \overline{\clQ}$ defines a closed geodesic in $\Sigma \setminus \overline{\clQ}$ freely homotopic to $\gamma$.

For any $\gamma_1$ and $\gamma_2$ that are either proper arcs or closed curves, denote by $I(\gamma_1, \gamma_2)$ their geometric intersection number, i.e. the minimal number of intersections of curves $\gamma'_1$ and $\gamma'_2$, where $\gamma'_1$ and $\gamma'_2$ are homotopic to $\gamma_1$ and $\gamma_2$, respectively. Similarly, we can define $I(\cdot, \cdot)$ for families of proper arcs resp.\ closed curves. We will below consider a proper arc $\tau$ with endpoints in $\overline{\clQ}$ and a collection of proper geodesic arcs $\clE$ such that for all $k\in \N$,
\begin{align}\label{intersection}
I(\varphi^{kn}(\tau), \clE) \geq (m-2)^k.
\end{align}
From this the Proposition will follow as in~\cite{FranksHandel1988}.



Let us introduce some terminology. Let $S_{\infty}$ the boundary at infinity of $\Hy$. Every geodesic in $\Hy$ has two endpoints in $S_{\infty}$. We say that a rectangle $Q=(\alpha_1, \alpha_2, \alpha_3, \alpha_4)$ in $\Hy \cup S_{\infty}$ with vertices in $S_{\infty}$ and adjacent edges $\alpha_1, \,\ldots,\, \alpha_4$ is a geodesic rectangle in $\Hy$ if $\alpha_1,\, \ldots, \,\alpha_4$ are geodesics. We keep the information of the ordering of the edges and say that the left vertical edge of $Q$ is $\alpha_1$, the right vertical edge of $Q$ is $\alpha_3$ and the horizontal edges are $\alpha_2$ and $\alpha_4$. We say that a geodesic $\beta$ in $\Hy$ intersects $Q$ horizontally from left to right if $\beta$ intersects each of its vertical edges, first the left and then the right vertical edge. We say that geodesic rectangle $Q$ intersects $Q'$ horizontally from left to right if both horizontal edges of $Q$, parametrized from the left vertical edges to the right vertical edges of $Q$, intersect $Q'$ horizontally from left to right. We say that $Q\prec Q'$, if there is a geodesic in $\Hy$ that first intersects $Q$ and then $Q'$ horizontally from left to right.

Let $j \in \{2,\,\ldots,\,m-1\}$. Writing $\clQ$ as in~\eqref{qqqq}, we let $e_j$ be a simple proper arc from $q_1^j$ to $q_3^j$ in $H_j \cap V_1$, $f_j$ a simple proper arc from $q_3^j$ to $q_5^j$ in $H_j \cap \varphi^n(H_1)$, $g_j$ a simple proper arc from $q_5^j$ to $q_7^j$ in $ H_j\cap V_m$, and $h_j$ a simple proper arc from $q_7^j$ to $q_1^j$ in $H_j \cap \varphi^n(H_m)$. The arcs $e_j, f_j, g_j$, resp.\ $h_j$, uniquely define homotopic simple geodesic arcs $\hat{e}_j, \hat{f}_j$, $\hat{g}_j$, resp.\ $\hat{h}_j$. Let $\hat{L}_j$ be the rectangle formed by $\hat{e}_j, \hat{f}_j$, $\hat{g}_j$, and $\hat{h}_j$. By the choice of $\clQ$, see~\eqref{1jm} and~\eqref{11mm}, the full rectangles $\hat{L}_j$ are up to their vertices completely contained in $\Sigma \setminus \overline{\clQ}$, and hence can be lifted to $\Hy$. Any lift $\widetilde{L}_j$ of $\hat{L}_j$ to $\Hy$ is a geodesic rectangle $Q=(\alpha_1,\, \ldots,\, \alpha_4)$ in $\Hy$, where $\alpha_1, \alpha_2, \alpha_3, \alpha_4$ are suitable lifts of $\hat{e}_j, \hat{f}_j$, $\hat{g}_j$, $\hat{h}_j$, respectively. Note that any two lifts of $\hat{L}_2,\,\ldots,\, \hat{L}_{m-1}$ to $\Hy$ are pairwise disjoint.


It is a theorem of Nielsen that any lift of $\varphi^n|_{\Sigma \setminus \clP}$ to $\Hy$ extends to a homeomorphism $F$ on $\Hy \cup S_{\infty}$, and let $F$ be such an extension. Let $Q$ be a geodesic rectangle in $\Hy$, then the image $F(Q)$ defines uniquely a geodesic rectangle $[F(Q)]$ in $\Hy$. Note that the endpoints of $\varphi^{n}(f_i)$ and $\varphi^{n}(h_i)$ for $i\in \{2, \,\ldots,\, m-1\}$ lie in $\overline{\clQ} \cap (H_1 \cup H_{m})$. It is now easy to verify that for all $j\in\{2, \,\ldots,\, m-1\}$ there is an arc $e'_j$ homotopic to $e_j$ and an arc $g'_j$ homotopic to $g_j$ such that for all $i\in \{2, \,\ldots,\, m-1\}$, $\varphi^n(f_i)$ intersects both $e'_j$ and $g'_j$ in exactly one point, and the order of those intersections coincides with the orientation of $f_i$ and $h_i$ if parametrized from $e_i$ to $g_i$. It follows that for any lift $Q$ of $\hat{L}_j$ there are lifts $Q_i$ of $\hat{L}_i$, $i=2,\, \ldots,\, m-1$, such that the geodesic rectangle $[F(Q)]$ intersects $Q_2, \,\ldots,\, Q_{m-1}$ all horizontally from left to right.

Let $q, q'$ be distinct points in $\overline{\clQ}$ and $\sigma$ a geodesic arc in $\Sigma \setminus \overline{\clQ}$ tracing a proper arc from $q$ to $q'$. Take a lift $\widetilde{\sigma}$ of $\sigma$ to $\Hy$. Let $D_j(\sigma)$ be the number of geodesic rectangles $Q=(\alpha_1,\, \ldots,\, \alpha_4)$ that arise as lifts of rectangles $\hat{L}_j$ and which $\sigma$ intersects horizontally from left to right. Let $D(\sigma) = \sum_{j=2}^{m-1} D_j(\sigma)$. These numbers obviously do not depend on the choice of lift of $\sigma$. Let $\sigma'$ be the geodesic arc in $\Sigma \setminus \clP$ that traces a proper arc homotopic to $\varphi^n(\sigma)$. We claim that
\begin{align}\label{Destimate}
D(\sigma') \geq (m-2)D(\sigma).
\end{align}
This can be seen as follows. Take a lift $\widetilde{\sigma}$ of $\sigma$ to $\Hy$. We can order the lifts of the $\hat{L}_j$, $j\in \{2,\, \ldots,\, m-1\}$ that $\widetilde{\sigma}$ intersects horizontally from left to right by $\prec$, i.e. these are geodesic rectangles $J_1,\, \ldots,\, J_k$, with $k=D(\sigma)$, and $J_i \prec J_j$ if $i<j$. Since $F|_{S_{\infty}}$ keeps the cyclic order of points at infinity of geodesics, it follows that the geodesic arc in $\Hy$ with the same endpoints as $F(\widetilde\sigma)$ has to intersect the geodesic rectangles $[F(J_1)]$, $[F(J_2)],\, \ldots,\, [F(J_k)]$ horizontally from left to right, and $[F(J_i)] \prec [F(J_j)]$ if $i< j$. Since each of the rectangles $[F(J_i)]$ intersects horizontally from left to right suitable lifts $Q_2, \,\ldots,\, Q_{m-1}$ of $\hat{L}_2,\, \ldots,\, \hat{L}_{m-1}$, \eqref{Destimate} follows.

Now choose $q$ and $q'$ two distinct points in $\overline{\clQ} \cap (H_1 \cup H_m)$ and $\tau$ an arc from $q$ to $q'$ with $D(\tau) = 1$. Furthermore, let $\clE = \bigcup_{j=2}^{m-1} \hat{e}_j$. From~\eqref{Destimate} the estimate~\eqref{intersection} follows, that is
\begin{align*}
I(\varphi^{kn}(\tau), \clE) \geq (m-2)^k.
\end{align*}
Choose now as in~\cite{FranksHandel1988} for each $k\in \N$ a closed curve $\gamma_k$ in $\Sigma\setminus \overline{\clQ}$ as the frontier of a small contour of the union of the geodesic representative of $\varphi^{kn}(\tau)$ and its two endpoints. For those $k \in \N$ for which $\varphi^{kn}(q), \varphi^{kn}(q') \in H_1 \cup H_m$, the curve $\gamma_k$ and $\clE$ have $2I(\varphi^{kn}(\tau), \clE)$ many intersections. Moreover there is no bigon formed by the union of $\clE$ with $\gamma_k$, and hence the number of those intersections is minimal among all curves homotopic to $\gamma_k$. Since $\varphi^{kn}(\gamma_0)$ is freely homotopic to $\gamma_k$ it follows that $\Gamma_{\pi_1}([\varphi^n, \clQ]) \geq \log(m-2)$.
\end{proof}

\begin{thebibliography}{99}
\bibitem[Kat84]{Katok1984} A. Katok, ``Nonuniform hyperbolicity and structure of smooth dynamical systems'', \emph{Proceedings of the International Congress of Mathematicians, Vol.\ 1, 2 (Warsaw, 1983)}, PWN, Warsaw, 1984, pp. 1245--1253.
\bibitem[KH95]{Hasselblatt-Katok} Anatole Katok and Boris Hasselblatt, \emph{Introduction to the modern theory of dynamical systems}, Encyclopedia of Mathematics and Its Applications, vol. 54, Cambridge University Press, 1995, with a supplementary chapter by Katok and Leonardo Mendoza.
\bibitem[FH88]{FranksHandel1988} John M. Franks and Michael Handel, ``Entropy and exponential growth of $\pi_1$ in dimension two'', \emph{Proc. Am. Math. Soc.} \textbf{102} (1988), no. 3, 753--760.
\bibitem[BP07]{BarreiraPesin} Luis Barreira and Yakov Pesin, \emph{Nonuniform hyperbolicity}, Encyclopedia of Mathematics and Its Applications, vol. 115, Cambridge University Press, 2007, Dynamics of systems with nonzero Lyapunov exponents.
\end{thebibliography}
\end{document}
